%% math environments

% definition style
\theoremstyle{definition}
\newtheorem{definition}{Definition}[section]
\newtheorem{remark}[definition]{Remark}
\newtheorem{example}[definition]{Example}
\newtheorem*{definition*}{Definition}
\newtheorem*{remark*}{Remark}
\newtheorem*{example*}{Example}

% plain style
\theoremstyle{plain}
\newtheorem{theorem}[definition]{Theorem}
\newtheorem{corollary}[definition]{Corollary}
\newtheorem{conjecture}[definition]{Conjecture}
\newtheorem{lemma}[definition]{Lemma}
\newtheorem{proposition}[definition]{Proposition}


%% enumeration items
\renewcommand{\labelenumi}{(\roman{enumi})}
\renewcommand{\labelenumii}{(\alph{enumii})}
\renewcommand{\labelenumiii}{(\roman{enumiii})}
\renewcommand{\labelenumiv}{(\roman{enumiv})}
\newcommand{\ima}{\operatorname{Im}}
\newcommand{\val}{\operatorname{val}}
\newcommand{\eps}{\varepsilon}
\newcommand{\dcup}{\operatorname{\dot{\cup}}}
\newcommand{\bbP}{\mathbb{P}}
\newcommand{\drawellipse}[5][]{
    \draw[#1] 
        let \p1 = (#2),
            \p2 = (#3),
            \n1 = {veclen(\x2-\x1, \y2-\y1)/2 + #4},
            \n2 = {#5},
            \n3 = {atan2(\y2-\y1, \x2-\x1)}
        in 
        [rotate=\n3] ($(#2)!0.5!(#3)$) ellipse (\n1 and \n2);
}